\section{Del stack de sistemas estable a la Great Transition}

Durante más de una década, el software de internet consolidó una vía de producción relativamente estable: el capital adquiere terrenos, electricidad y centros de datos; los chips llegan a las salas de servidores; las plataformas en la nube empaquetan el hardware como recursos invocables; los equipos entrenan modelos y construyen aplicaciones sobre ellas, y después los mecanismos de seguridad, cumplimiento normativo y gobernanza llevan el sistema al mundo real. La IA generativa no se limitó a añadir una categoría de aplicaciones en la cima de este stack; al mismo tiempo elevó la presión sobre las capas de consumo energético, interconexión, chips, datos, modelos, productos y gobernanza, lo que obliga a renegociar los antiguos límites de responsabilidad y estructuras de precios.

\begin{figure}[H]
\centering
\includegraphics[width=0.94\textwidth]{images/00-10-00-course-in-a-nutshell.jpg}
\caption{El curso organiza el contenido de diez semanas con una cadena de todo el stack que va del capital, la energía y los centros de datos a los modelos, las aplicaciones y la gobernanza.}
\figsource{Video oficial de Stanford Online 00:10:00, `images/00-10-00-course-in-a-nutshell.jpg`.}
\end{figure}

\begin{knowledgebox}{Lectura de la figura: primero en la dirección de las dependencias, después en la dirección del valor}
De abajo hacia arriba, el capital, los terrenos, la electricidad, los chips y los centros de datos aportan las condiciones físicas para el cómputo; la cloud/software infrastructure convierte el hardware heterogéneo en recursos orquestables; los modelos y los agent consumen esos recursos para producir capacidades; las aplicaciones incorporan esas capacidades a los flujos de trabajo, y la capa de gobernanza se ocupa de la seguridad, la responsabilidad y la aceptación social. De arriba hacia abajo, los ingresos de las aplicaciones y el valor público determinan, a su vez, en qué tipo de capacidad de cómputo, modelos e instalaciones está dispuesto el capital a seguir invirtiendo. Ambas direcciones coexisten, por lo que cualquier «avance» en una sola capa puede quedar frenado por otra.
\end{knowledgebox}

\begin{table}[H]
\centering
\begin{tabular}{L{0.18\textwidth}L{0.34\textwidth}L{0.38\textwidth}}
\toprule
Capa & Recursos o interfaces principales & Preguntas de auditoría de Frontier Systems \\
\midrule
Capital y energía & Financiamiento, terrenos, red eléctrica, compra de electricidad a largo plazo & ¿El ciclo de inversión puede ajustarse a la incertidumbre de la tecnología y de la demanda? \\
Centros de datos y chips & Salas de servidores, aceleradores, redes, enfriamiento & ¿Qué hardware es realmente intercambiable y dónde hay dependencia de proveedor? \\
Nube e infraestructura de software & Planificación, virtualización, almacenamiento, observabilidad & ¿Cómo se convierten los recursos heterogéneos en interfaces de servicio estables? \\
Modelos y agent & Pesos, datos, entornos, retroalimentación & ¿La capacidad proviene de parámetros estáticos o de un ciclo cerrado continuo? \\
Aplicaciones y despliegue & Usuarios, flujos de trabajo, distribución, ingresos & ¿Cómo obtiene el producto context de alta calidad y señales de verificación? \\
Seguridad y gobernanza & Normas, auditoría, responsabilidad, instituciones & ¿Quién define los límites y quién asume el costo de los fallos? \\
\bottomrule
\end{tabular}
\caption{El stack del curso no es una lista de empresas, sino una tabla de dependencias de recursos y de asignación de responsabilidades.}
\end{table}

\subsection{Declaración de intereses del ponente y posición histórica}

Para valorar la solidez de la evidencia de un curso sobre la industria, esta sección vuelve a situar en el panorama la línea de tiempo de participación del ponente. Las inversiones tempranas, la participación en consejos de administración, las colaboraciones en productos o la fundación de empresas aportan detalles difíciles de obtener en materiales públicos, pero también pueden hacer que ciertos temas reciban más atención. Un buen método de estudio no consiste en descartar todas las opiniones porque exista un conflicto de intereses, sino en distinguir entre observación de primera mano, narrativa corporativa, inferencia macro y juicio de valor, y buscar evidencia verificable para cada una.

\begin{figure}[H]
\centering
\includegraphics[width=0.94\textwidth]{images/00-16-45-background-disclosures.jpg}
\caption{Declaración del ponente, en forma de línea de tiempo, de sus inversiones y de los proyectos que fundó o en los que participó.}
\figsource{Video oficial de Stanford Online 00:16:45, `images/00-16-45-background-disclosures.jpg`.}
\end{figure}

\begin{warningbox}{Lectura de la figura: la declaración no es un descargo de responsabilidad decorativo}
La línea de tiempo indica al lector qué juicios pueden provenir de experiencia interna y cuáles solo pueden ser analogías externas. Por ejemplo, alguien que lleva mucho tiempo involucrado en empresas de modelos y en el mercado de capacidad de cómputo percibe con más facilidad la capital allocation y el deployment feedback, pero no necesariamente representa a todos los equipos de investigación, las instituciones públicas o los desarrolladores pequeños y medianos. Cuando más adelante se citen las opiniones del ponente sobre las context wars, los precios de las GPU y los arreglos institucionales, deben marcarse como juicios basados en la experiencia con una fuente identificada, y no como leyes objetivas sin sujeto.
\end{warningbox}

\subsection{Cuatro tipos de restricciones de capacidad}

Con el mapa de todo el stack, el curso agrupa los límites del progreso de la IA en cuatro grupos: context, compute, capital y culture. Este orden merece atención. El debate público suele tratar el compute como el único bien escaso, pero el cómputo solo se convierte en capacidad cuando entra en contacto con las tareas, los entornos y la retroalimentación adecuados; el capital determina la velocidad de expansión y la asunción de riesgos, y la cultura influye en qué está dispuesto a probar un equipo, qué fallos reconoce y qué valores incorpora a los productos y a las instituciones.

\begin{figure}[H]
\centering
\includegraphics[width=0.9\textwidth]{images/00-17-45-capability-bottlenecks.jpg}
\caption{Los cuatro tipos de restricciones al progreso de las capacidades de la IA: context, compute, capital y culture.}
\figsource{Video oficial de Stanford Online 00:17:45, `images/00-17-45-capability-bottlenecks.jpg`.}
\end{figure}

\begin{importantbox}{Cómo interactúan los cuatro tipos de restricciones}
\term{Context} determina qué tareas, herramientas y retroalimentación ve el sistema; \term{compute} aporta el presupuesto de entrenamiento e inferencia; \term{capital} convierte las expectativas de ingresos futuros en inversión en salas de servidores, chips y talento; \term{culture} determina cómo la organización comparte los errores, gestiona la seguridad, elige objetivos y coordina la colaboración. Cuando una capa es suficiente de momento, la presión aparece en las demás. Al diagnosticar un sistema hay que buscar la restricción marginal actual; no se puede reutilizar siempre la misma respuesta.
\end{importantbox}

\begin{knowledgebox}{Términos clave: bottleneck no es una etiqueta permanente}
Aquí, bottleneck designa el eslabón que más limita la producción total cuando se agrega una unidad de recurso en el momento actual. Cambia con la escala, la tarea, la región y el tiempo. Por ejemplo, a un equipo pequeño puede faltarle primero capital; a un gran laboratorio, primero electricidad e interconexión, y a un coding agent maduro, sobre todo entornos ejecutables y retroalimentación de usuarios reales. Si el curso cruza capas con frecuencia más adelante, es precisamente porque las restricciones se desplazan y también se amplifican entre sí.
\end{knowledgebox}

\subsection{Resumen de la sección}

La Great Transition significa que todo el stack técnico-económico está cambiando a la vez. El mapa de todo el stack aporta las relaciones de dependencia, la línea de tiempo de la declaración de intereses ayuda a juzgar la procedencia de las opiniones y los cuatro tipos de restricciones ofrecen un marco de diagnóstico dinámico. A continuación nos centramos en la primera restricción: cómo el context forma, mediante entornos, retroalimentación y aprendizaje por refuerzo, un ciclo cerrado de capacidades sostenible.
